\documentclass[10pt]{article}

\usepackage[T1]{fontenc}
\usepackage{lmodern}
\usepackage[margin=0.62in,top=0.52in,bottom=0.55in]{geometry}
\usepackage{amsmath,amssymb}
\usepackage{graphicx}
\usepackage{microtype}
\usepackage{xcolor}
\usepackage[hidelinks]{hyperref}

\definecolor{MSSNavy}{HTML}{12304A}
\definecolor{MSSOrange}{HTML}{C65D16}
\definecolor{MSSGreen}{HTML}{2F855A}
\definecolor{MSSGray}{HTML}{627D98}

\newcommand{\compactheading}[1]{%
  \vspace{0.32em}\noindent{\color{MSSOrange}\bfseries\large #1}\par\vspace{0.16em}}
\newcommand{\resulttag}[1]{\textcolor{MSSGreen}{\textbf{#1}}}

\setlength{\parindent}{0pt}
\setlength{\parskip}{0.32em}

\begin{document}

{\color{MSSOrange}\bfseries\small REFEREE REPORT}\hfill
{\color{MSSGray}\small 15 August 2026}\par
{\color{MSSNavy}\bfseries\LARGE Referee Report: The Saving Glut of the Rich}\par
{\color{MSSGray}\small Review of Mian, Straub, and Sufi (July 2025)}
\vspace{0.25em}\hrule height 1.1pt\vspace{0.35em}

\resulttag{Recommendation: revise and resubmit.}
This is an ambitious, original, and highly relevant paper. Its central
innovation is to combine Financial Accounts data with a Leontief-style
operator that traces claims through intermediary cross-holdings to their
ultimate owners and uses. This reveals a major within-household shift that
aggregate national accounts conceal: the rise in top-1-percent saving and its
counterpart in household and government debt. The treatment of valuation
changes, retained corporate earnings, and national-accounts closure is
particularly valuable. Our reconstruction using the older supplied inputs
reproduces the main patterns closely, increasing confidence that the paper
captures a real and economically important fact. I recommend publication after
a focused revision addressing the following measurement and interpretation
issues.

\compactheading{1. Show saving rates over time and by inflation-adjusted wealth levels}

Figure 6 compares two long-period averages by wealth percentile. This is a
useful summary, but percentile rank is not the only relevant horizontal axis.
If households become richer over time and consumption rises less than
proportionally with wealth, an increasing saving rate at high ranks could
partly reflect movement along a stable, non-homothetic saving schedule rather
than a change in behavior conditional on resources. The paper should therefore
show the same saving-rate definition against inflation-adjusted wealth levels,
ideally using fixed real-wealth bins in addition to percentile bins.

Our diagnostic re-plots the same percentile-cohort saving rates at their mean
net wealth per adult in 2018 dollars (Referee Figure 1). The result does
\textbf{not} make the post-1982 divergence disappear. The pre-1982
90th-percentile cohort has mean wealth of about \$303,000 and a 31.5 percent
saving rate, while the post-1982 80th-percentile cohort has similar mean wealth
of about \$286,000 but saves only 7.4 percent. The top 1 percent moves from
about \$3.0 million to \$6.9 million in mean wealth while its saving rate rises
from 28.8 to 36.2 percent. These comparisons suggest both a rightward shift in
wealth and a change in the saving-rate profile. They remain descriptive because
cohort means are not fixed-dollar bins and the data are repeated
cross-sections; this is precisely why the authors should perform the
fixed-real-wealth-bin exercise on their exact 1963--2019 microdata.

The two-period average also hides considerable time variation (Referee Figure
2). In our older-input reconstruction, the five-year top-1 rate is 35.6 percent
in the window ending in 1982, peaks at 68.5 percent in the window ending in
2008, and falls to 24.1 percent by 2016. The bottom-40 rate moves from -3.3 to
-15.3 and back to -3.2 percent over the same windows. This is consistent with
strong crisis-era movements at both ends, in opposite directions, but it is
not a causal Great Recession estimate: overlapping windows blur timing and the
2008 window contains four pre-crisis years. The paper should add annual or
non-overlapping windows, recession shading, and a decomposition into financial
assets, housing, debt write-downs, and the equity-valuation residual.

A standard logarithmic color scale is invalid because saving rates can be zero
or negative. I suggest a linear, zero-centered main heatmap and a clearly
labelled symmetric-logarithmic or inverse-hyperbolic-sine diagnostic,
accompanied by selected-percentile lines. If the intended object is
\emph{change}, the heatmap should plot deviations from each percentile's
pre-1982 mean rather than levels.

\compactheading{2. Bound what the unveiling and saving assumptions identify}

The proportional allocation rule is transparent, but claims issued by an
intermediary differ in seniority, maturity, collateral, guarantees, and risk.
The operator therefore identifies proportional ultimate claim ownership, not
necessarily loss incidence, control, or causation. The paper should report the
observed, balanced, and imputed shares of the direct matrix; the residual
share; and the sensitivity of Figures 8--10 to perturbations of the largest
imputed links and alternative instrument-specific allocations.

Likewise, NIPA closure is an essential accounting check but does not identify
how the residual equity valuation term is distributed. Because equity
ownership is concentrated, plausible alternatives for equity inflation,
corporate saving, housing returns, and debt write-downs should be propagated
into uncertainty or sensitivity bands for Figures 5 and 6. These diagnostics
are especially important around 2008. Finally, causal verbs such as
``fueled'' or ``drove'' should be reserved for identified mechanisms; the
baseline establishes that saving is ultimately \emph{absorbed by} or
\emph{held against} particular debt claims.

\compactheading{3. Provide a version-matched replication package}

The February 2021 package reproduces its own paper version, but it generally
ends in 2016 and uses the earlier round-by-round unveiling. The revised paper
should archive the exact 2025 inputs, completed annual matrices, crosswalks,
valuation construction, and figure scripts behind the 2019 results. This is
particularly important for a paper whose methodology is itself a major
contribution.

Overall, these requests would sharpen rather than overturn the main finding.
The paper offers a useful framework and compelling evidence that rising
inequality has changed the composition and ultimate use of saving.

\clearpage
{\color{MSSNavy}\bfseries\Large Referee Figure 1 --- Saving rates by rank and real wealth}\par
\vspace{0.25em}
\begin{center}
  \includegraphics[width=0.94\textwidth]{../Results_Proposal/figures/figure6_percentile_vs_wealth_level.png}
\end{center}
\vspace{-0.35em}
{\small\textbf{Finding.} Re-indexing the same cohorts by mean real wealth does
not remove the post-1982 divergence. The comparison is descriptive and does
not replace the proposed fixed-real-wealth-bin estimate.}\par
{\color{MSSGray}\scriptsize Source: generated by the independent Figure 6
reconstruction from February 2021 input-vintage data through 2016. Net wealth
is measured per adult in 2018 dollars.}

\vspace{0.55em}
{\color{MSSNavy}\bfseries\Large Referee Figure 2 --- Saving rates over time}\par
\vspace{0.2em}
\begin{center}
  \includegraphics[width=0.80\textwidth]{../Results_Proposal/figures/figure6_rolling5_heatmap.png}
\end{center}
\vspace{-0.35em}
{\small\textbf{Finding.} Five-year saving rates vary substantially over time
and across the wealth distribution, including sharp crisis-era movements near
both tails. Overlapping windows do not identify a causal recession effect.}\par
{\color{MSSGray}\scriptsize Source: generated by the independent Figure 6
reconstruction from February 2021 input-vintage data through 2016. Colors show
trailing five-year mean net saving rates.}

\end{document}
